\subsection{A window test using only the fixed target and ordinary primes}\label{pref:section}

Fix a primitive character \(\psi\) of order \(o\). Write
\(\chi_a=(\psi^a)^\ast\) for the primitive inducers in
Proposition~\ref{ar:mean}, with the same weights \(\omega_a\). Put
\[
 A(n)=\sum_{a=0}^{o-1}\omega_a\Re\chi_a(n),\qquad
 c_\psi=1+\gamma+\sum_{a=1}^{o-1}\omega_a
                      \Re\frac{L'(1,\chi_a)}{L(1,\chi_a)} .
\]
Then \(2/3\le A(n)\le4/3\). At a prime unramified in the conductor
of \(\psi\), character orthogonality gives
\[
 A(p^j)=\frac{3o+o\,\mathbf1_{\psi(p)^j=1}}{3o+1}.
\]
At the other primes the definition by individual primitive inducers
applies. By \eqref{eq:centred}, the fixed retained function is
\[
 F(x)=G_{\mathcal F_\psi}(x)
     =\sum_{n\le x}A(n)\Lambda(n)(x/n-1)-x\log x+c_\psi x .
\]

\begin{proposition}[Removing the higher powers on a finite window]\label{pref:powers}
Let \(2\le X\le x\le Z\), and define
\[
\begin{aligned}
 b_Z&=\sum_{p\le Z}\log p\sum_a\omega_a\Re
           \frac{\chi_a(p)^2}{p(p-\chi_a(p))},\\
 W_Z(x)&=\sum_{p\le x} A(p)\log p\,(x/p-1)
                       -x\log x+(c_\psi+b_Z)x .
\end{aligned}
\]
The difference is nonnegative, with the exact formula
\[
 W_Z(x)-F(x)=\sum_{p\le Z}\sum_{j\ge2}
           A(p^j)\log p\min(1,x/p^j).
\]
It is \(O(\sqrt x)\) with an absolute constant. Removing the displayed
\(j=2\) term leaves an absolute \(O(x^{1/3})\). For each fixed
\(\psi\), uniformly for \(Z\ge x\to\infty\),
\[
 W_Z(x)-F(x)=
       2\,\frac{3o+\gcd(2,o)}{3o+1}\sqrt{x}+o_\psi(\sqrt x).
\]
\end{proposition}
\begin{proof}
For each fixed prime, summing \(\chi_a(p)^j/p^j\) for \(j\ge2\)
gives the geometric fraction in \(b_Z\), including when \(\chi_a(p)=0\).
Subtract the higher ramps from \(xb_Z\) and use
\(u-(u-1)_+=\min(1,u)\) for \(u\ge0\). For \(r=2,3\) the logarithmic
mass of powers \(p^j\le t\), \(j\ge r\), is \(O(t^{1/r})\).
Indeed the terms \(j=r\) satisfy Chebyshev's bound and the rest are
\(O(t^{1/(r+1)}\log(2t))\). Integrate this bound using
\(\min(1,x/n)=x\int_x^\infty\mathbf1_{n\le t}\,dt/t^2\).

On the reduced classes of the conductor, the fraction with
\(\psi^2=1\) is \(\gcd(2,o)/o\). The prime number theorem in fixed
arithmetic progressions (for example
\href{https://arxiv.org/abs/1802.00085}{Bennett, Martin, O'Bryant and
Rechnitzer}, Theorem 1.2 for moduli at least three) therefore gives
\(\sum_{p\le v}A(p^2)\log p=(3o+\gcd(2,o))v/(3o+1)+o_\psi(v)\).
Integration with \(v=\sqrt t\) gives the asserted leading term for
all squares. Restricting the primes to \(p\le Z\) changes that
integral by \(O(x/Z)\), since
\(\sum_{p>Z}\log p/p^2\ll 1/Z\) by partial summation.
\end{proof}

\begin{proposition}[Finite test for every real point of the window]\label{pref:test}
For the empty prefix and each prime \(p\le Z\), put
\[
 t_p=\sum_{\ell\le p}A(\ell)\log\ell,\qquad
 m_p=\sum_{\ell\le p}A(\ell)\log\ell/\ell ,
\]
where \(\ell\) runs over primes and the empty sums vanish.
Let \(0<\alpha<1\), \(k\ge0\). For each prefix choose a minimizer
\(z_p\in[X,Z]\) of \(\log z+t_p/z+kz^{\alpha-1}\). Then
\[
 W_Z(x)\le k x^\alpha\quad(X\le x\le Z)
\]
holds if and only if every prefix satisfies
\[
       m_p+c_\psi+b_Z\le\log z_p+t_p/z_p+kz_p^{\alpha-1}.
\]
For fixed \(k,\alpha\) and large \(X\), it suffices to check \(X,Z\)
and just the roots
\[
 z-k(1-\alpha)z^\alpha=t_p
\]
lying between the last prime of the active prefix and the next prime
inside the window. An end gap is intersected with \([X,Z]\).
\end{proposition}
\begin{proof}
Positivity of \(A\) gives an envelope:
\(\sum_{p\le x}A(p)\log p(x/p-1)=\max_p(xm_p-t_p)\).
At \(x\) the prefix through its last prime attains equality. Divide by
\(x\) and require each line to lie below the proposed bound, then
minimize its right side over the whole interval. For \(t_p>0\) the
identity
\[
 m_p+c_\psi+b_Z-\log x-t_p/x
 =m_p+c_\psi+b_Z-1-\log t_p
                  -\bigl(t_p/x-1-\log(t_p/x)\bigr)
\]
also gives the exact loss away from the maximizer of the normalized
line at \(x=t_p\).

For the shorter list differentiate \((W_Z(x)-kx^\alpha)/x\) on
each prime gap. Its derivative is
\((t_p-x+k(1-\alpha)x^\alpha)/x^2\).
The function \(x-k(1-\alpha)x^\alpha\) is strictly increasing for
large \(x\). At a prime the derivative jumps upwards, so an interior
maximum must occur at a gap root. There the test simplifies to
\(m_p+c_\psi+b_Z-1-\log z\le k\alpha z^{\alpha-1}\).
The empty prefix is covered by the same derivative formula.
For \(\alpha=1/2\) the root is explicit:
\(\sqrt z=k/4+\sqrt{t_p+k^2/16}\).
The argument applies verbatim to all prime powers with masses
\(A(p^j)\log p\) and centre \(c_\psi\), giving the direct test for \(F\).
\end{proof}

\paragraph{The implication for the actual mean.}
Keep the common character, scale and size hypotheses of
Corollary~\ref{ar:meanhost}, including its discarded-twist bounds.
Write \(t=n_T,\ B=tmb,\ \delta=n_2/n\), and let \(E_x\) be the
right side of \eqref{scope:missing}. By dropping two nonnegative
terms and using \(a_E\ge xb\), we have
\[
 \frac{2E_x}{B\delta}
 \ge \frac{2x^{5\varepsilon}}{\delta}
       \left(\frac{x}{m}+2\frac{n}{t}m x^\varepsilon\right)
 \ge \frac{4\sqrt2}{\delta}\sqrt{\frac nt}\,
                   x^{1/2+11\varepsilon/2}.
\]
Consequently, a bound \(W_Z(x)\le x^{1/2+11\varepsilon/2}\)
on arbitrarily late common windows \(X=Y,\ Z=Y^{3/2}\) suffices
for \eqref{scope:average}. Indeed \(t\le n\), \(x^{-\varepsilon}
\le\delta\le1\), \(F\le W_Z\), and
Corollary~\ref{ar:meanhost} bounds the normalized mean minus \(F\)
by \(O_\psi(x^{1/2+\varepsilon}+(1+x^\varepsilon)\log(2x)
+1/(x\delta))\) for the fixed uniform host constants.
Conversely, \eqref{scope:average} on these data implies
\(W_Z(x)\le 2E_x/(B\delta)+O_\psi(x^{1/2+\varepsilon})\).
The finite test only uses the primes at most \(Z\) and a constant
of the fixed target. The original host assumptions enter at the
comparison with the mean.

\begin{corollary}[Size of the required prime bound on the actual ranks]\label{pref:scale}
On the full rank page of Proposition~\ref{scope:reselection} put
\[
     S_x=\delta^{-1}\sqrt{x}\,x^{97\varepsilon/16}.
\]
Then \(2S_x\le 2E_x/(B\delta)\le28S_x\) at every sufficiently
large scale. In particular \(W_Z\le S_x\) on a common window
implies \eqref{scope:average} on its data, while
\eqref{scope:average} there implies \(W_Z\le(28+o(1))S_x\).
If the reselected sets can be realized as in
Corollary~\ref{scope:reselection-loss}, a fixed implicit constant
in \(W_Z\ll S_x\) suffices for both inequalities on the new data.
\end{corollary}
\begin{proof}
Write \(y=x^\varepsilon\), \(v=\sqrt{x}y^{17/16}\).
The page gives \(\sqrt{x}y^{-1/16}\le m\le2\sqrt{x}y^{-1/16}\),
\(y^{1/8}/2\le n/t\le2y^{1/8}\), \(a_E\le2xb\),
\(xm\le k\), and \(ky^2\le n^2\sqrt{x}\) eventually.
Divide the four terms in the parentheses of \eqref{scope:missing}
by \(B\). The first is at most \(2x/m\le2v\);
the second is at most \(2xm/k\le2v\).
The third is \(2(n/t)my\), between \(v\) and \(8v\).
The fourth is at most \((n/t)\sqrt{x}\le2v\) since \(n_2\le n\).
Multiplication by \(2y^5/\delta\) proves the budget bounds.
The discarded terms, the higher powers and the normalized mean
remainder together cost \(o(S_x)\), uniformly at the fixed
\(\psi,\varepsilon\). The final assertion follows by taking a
constant \(g\) in Corollary~\ref{scope:reselection-loss}.
\end{proof}

\paragraph{The arithmetic estimate left by the test.}
Writing \(\vartheta_A(v)=\sum_{p\le v}A(p)\log p\) and
\(b_\infty=\lim b_Z\), there is an exact tail identity
\[
 \frac{W_Z(x)}x
 =\int_x^\infty\frac{v-\vartheta_A(v)}{v^2}\,dv-(b_\infty-b_Z).
\]
To see this, the prime number theorem with the fixed-modulus
de la Vall\'ee Poussin remainder gives
\(\sum_{n\le v}A(n)\Lambda(n)/n=\log v+C+o_\psi(1)\)
by partial summation. To identify \(C\), write the corresponding
Dirichlet series at \(1+s\) as the integral of the last partial sum
against \(s v^{-1-s}dv\) on \([1,\infty)\).
As \(s\downarrow0\) the integral is \(1/s+C+o(1)\);
the logarithmic derivatives of the primitive inducers give
\(1/s+1-c_\psi+O_\psi(s)\). Removing the absolutely summable
higher powers therefore gives
\(m_p=\log p+1-c_\psi-b_\infty+o_\psi(1)\).
Together with \(\vartheta_A(x)/x\to1\), this shows
\(W_\infty(x)/x\to0\). Integrate its derivative
\((\vartheta_A(x)-x)/x^2\) and restore the finite calibration.
The qualitative prime number theorem identifies the constant and
the known square correction. A one-sided bound on this signed
tail at the required power throughout arbitrarily late windows is
the outstanding estimate.

\begin{corollary}[Strength of the prime test]\label{pref:zeros}
For a fixed \(\psi\), the following are equivalent.
\begin{enumerate}
\item For every \(0<\eta<1/2\) there are arbitrarily late windows
\([Y,Z]\), \(Z=Y^{3/2}\), on which \(W_Z(x)\le x^{1/2+\eta}\)
at every real \(x\).
\item All the retained \(L(s,\chi_a)\) are zero-free in
\(\Re s>1/2\), away from the principal pole.
\item For every \(0<\eta<1/2\), the same inequality holds for all
sufficiently large \(x\) and every \(Z\ge x\).
\end{enumerate}
\end{corollary}
\begin{proof}
In the first condition \(F\le W_Z\). Theorem~\ref{an:zero} applied
at \(\theta=1/2+\eta\) excludes zeros to the right of that line;
then let \(\eta\) decrease to zero. For the converse use the
absolutely convergent inverse Mellin integral for the uncentred ramp
on \(\Re s=2\), with integrand
\[
        -\frac{L'}{L}(s,\chi_a)\,\frac{x^s}{s(s-1)} .
\]
Under the second condition move the contour to
\(\sigma=1/2+\eta/2\). At large heights, the fixed-character
logarithmic derivative is \(O_{\psi,\eta}(\log(2+|\Im s|))\)
uniformly between these lines. Indeed its local partial fraction
formula involves the zeros within one unit of height, of which
there are \(O_\psi(\log(2+|\Im s|))\); each denominator has modulus
at least \(\eta/2\). The horizontal integrals tend to zero.
The only crossed residue is at one. It gives
\(-x(L'/L)(1,\chi_a)\) for a nonprincipal character and
\(x(\log x-1-\gamma)\) for the principal one.
The centres cancel these terms exactly. The new vertical integral
is \(O_{\psi,\eta}(x^{1/2+\eta/2})\) because its kernel decays
quadratically in height. Proposition~\ref{pref:powers} now yields
the third condition, uniformly in \(Z\). This implies the first.
\end{proof}
For the trivial target the second condition is equivalent to RH by
the functional equation. Thus proving the finite test on arbitrarily
late windows for every sufficiently small exponent would itself give
a proof of RH for this target.

\paragraph{A complete window with a ramified quartic target.}
Take \(\psi=\chi_{5,4}\chi_{3,2}\), where
\(\chi_{5,4}(2)=i\). Its square is induced by the quadratic
character of conductor five. Thus \(13A(3^j)=13+(-1)^j\),
whereas \(A(5^j)=1\). On the complete window \(X=10^4\),
\(Z=X^{3/2}=10^6\), outward interval arithmetic gives
\[
              W_Z(x)<3\sqrt{x},\qquad F(x)<\tfrac14\sqrt{x}
                    \qquad(X\le x\le Z).
\]
The two computations use respectively the first powers only and all
prime powers. There are \(78{,}498\) primes and \(78{,}734\) prime
powers. The upward jumps of the normalized derivative leave only
\(1{,}528\) and \(2{,}713\) possible interior maxima, respectively.
Both endpoints are also tested. The resulting upper bounds for
\((W_Z-3\sqrt{x})/x\) and \((F-\sqrt{x}/4)/x\) are less than
\(-0.00067335\) and \(-0.00022959\).

Here is how the constants and tests can be reproduced. For
\(p\nmid15\) with \(d=\operatorname{ord}(\psi(p))\), the local
coefficient of \(\log p\) in \(13b_Z\) is
\[
    \frac{12}{p(p-1)}+
       4\left(\frac1{p^d-1}-\frac{\mathbf1_{d=1}}p\right).
\]
For \(p=3\) it is \(13/(3\cdot2)+1/(3\cdot4)\), and for \(p=5\)
it is \(13/(5\cdot4)\). The resulting enclosures are
\(0.7536664167<b_Z<0.7536664168\) and
\(1.7030116351<c_\psi<1.7030116379\).
To enclose a real or imaginary component of \(L(1)\) or \(-L'(1)\),
sum its zero-mean signed period against \(f(t)=1/t\) or
\(f(t)=\log(t)/t\) to a whole period ending at \(m\ge10^5\).
Let \(P\) be its sequence of partial sums over a period, \(\bar P\)
its period mean, and \(D\) the maximum absolute partial sum of
\(P-\bar P\). Twice summing by parts gives a remaining tail
\(\bar P f(m+1)\) with error at most
\(D(f(m+1)-f(m+2))\). Both first differences are positive and
decreasing here. The script combines these intervals by complex
division and encloses the candidate roots by integer square roots.
The source is \texttt{verification/prime\_prefix\_window.py}.
This calculation verifies the fixed-character functions on one
complete window. The common-host comparison has an asymptotic
threshold, and arbitrarily late windows remain necessary.
